
% TFG.tex
% TFG
% Autor
% Sevilla, 16 de julio de 2016
% =============================================================================

\documentclass[a4paper,12pt,twoside]{book}

%%%%%%%%%%%%%%%%%%%%%%%%%%%%%%%%%%%%%%%%%%%%%%%%%%%%%%%%%%%%%%%%%%%%%%%%%%%%%%
%%  Paquetes adicionales                                                   %%
%%%%%%%%%%%%%%%%%%%%%%%%%%%%%%%%%%%%%%%%%%%%%%%%%%%%%%%%%%%%%%%%%%%%%%%%%%%%%%

\usepackage[utf8x]{inputenc}       % Acentos de UTF8
\usepackage[spanish]{babel}        % Castellanización.
\usepackage[T1]{fontenc}           % Codificación T1 con European Computer
                                   % Modern.  
\usepackage{graphicx}
\usepackage{fancyvrb}              % Verbatim extendido
\usepackage{makeidx}               % Índice
\usepackage{amsmath}               % AMS LaTeX
\usepackage{amsthm} 
\usepackage{latexsym}
\usepackage[colorinlistoftodos
           , backgroundcolor 
= yellow
           , textwidth = 4cm
, shadow
           , spanish]{todonotes}
% Fuentes
\usepackage{mathpazo}              % Fuentes semejante a palatino
\usepackage[scaled=.90]{helvet}
\usepackage{cmtt}
\renewcommand{\ttdefault}{cmtt}
\usepackage{a4wide}
% \usepackage{xmpincl}               % Incluye metadato de licencia CC.

% Tikz
\usepackage{tkz-berge}
\usetikzlibrary{shapes,trees}
%añadidos por mi:
\usepackage{parskip}
\usepackage{cite} % para contraer referencias

\newtheorem{theorem}{Teorema}
\newtheorem{defi}{Definición}
\newtheorem{prop}{Proposición}
\newtheorem{dem}{Prueba}

% Desactivar <,> cuando se hacen dibujos con tikz.
\tikzset{
every picture/.append style={
  execute at begin picture={\deactivatequoting},
  execute at end picture={\activatequoting}
  }
}

% Márgenes
\usepackage[margin=1in]{geometry}

% Control de espacios en la tabla de contenidos:
\usepackage[titles]{tocloft}
\setlength{\cftbeforechapskip}{2ex}
\setlength{\cftbeforesecskip}{0.5ex}
\setlength{\cftsecnumwidth}{12mm}
\setlength{\cftsubsecindent}{18mm}

% Control de listas
% Elimina espacio entre item y párrafo y coloca item en el margen izquierdo
% \usepackage{enumitem}
% \setlist[enumerate,itemize]{noitemsep, nolistsep, leftmargin=*}

\usepackage{minitoc}

% Doble espacio entre líneas
\usepackage{setspace}
\onehalfspacing


% \linespread{1.05}                % Distancia entre líneas
\setlength{\parindent}{5mm}        % Indentación de comienzo de párrafo
\deactivatetilden                  % Elima uso de ~ para la eñe
\raggedbottom                      % No ajusta los espacios verticales.

\usepackage[%
  pdftex,
  pdfauthor={Eduardo Paluzo},%
  pdftitle={LPO en Haskell},%
  pdfstartview=FitH,%
  bookmarks=false,%
  colorlinks=true,%
  urlcolor=blue,%
  unicode=true]{hyperref}      

\usepackage{tikz}

%%%%%%%%%%%%%%%%%%%%%%%%%%%%%%%%%%%%%%%%%%%%%%%%%%%%%%%%%%%%%%%%%%%%%%%%%%%%%%
%%  Cabeceras                                                              %%
%%%%%%%%%%%%%%%%%%%%%%%%%%%%%%%%%%%%%%%%%%%%%%%%%%%%%%%%%%%%%%%%%%%%%%%%%%%%%%

\usepackage{fancyhdr}

\addtolength{\headheight}{\baselineskip}

\pagestyle{fancy}

\cfoot{}

\fancyhead{}
\fancyhead[RE]{\bfseries \nouppercase{\leftmark}}
\fancyhead[LO]{\bfseries \nouppercase{\rightmark}}
\fancyhead[LE,RO]{\bfseries \thepage}

%%%%%%%%%%%%%%%%%%%%%%%%%%%%%%%%%%%%%%%%%%%%%%%%%%%%%%%%%%%%%%%%%%%%%%%%%%%%%%
%%  Definiciones                                                           %%
%%%%%%%%%%%%%%%%%%%%%%%%%%%%%%%%%%%%%%%%%%%%%%%%%%%%%%%%%%%%%%%%%%%%%%%%%%%%%%

\input definiciones
\def\mtctitle{Contenido}

%%%%%%%%%%%%%%%%%%%%%%%%%%%%%%%%%%%%%%%%%%%%%%%%%%%%%%%%%%%%%%%%%%%%%%%%%%%%%%
%%  Título                                                                 %%
%%%%%%%%%%%%%%%%%%%%%%%%%%%%%%%%%%%%%%%%%%%%%%%%%%%%%%%%%%%%%%%%%%%%%%%%%%%%%%

\title{\Huge Título del TFG}
\author{Autor}
\date{\vfill \hrule \vspace*{2mm}
  \begin{tabular}{l}
      \href{http://www.cs.us.es/glc}
           {Grupo de Lógica Computacional} \\
      \href{http://www.cs.us.es}
           {Dpto. de Ciencias de la Computación e Inteligencia Artificial} \\
      \href{http://www.us.es}
           {Universidad de Sevilla}  \\
      Sevilla, 16 de junio de 2016 (Versión de \today)
  \end{tabular}\hfill\mbox{}}




%%%%%%%%%%%%%%%%%%%%%%%%%%%%%%%%%%%%%%%%%%%%%%%%%%%%%%%%%%%%%%%%%%%%%%%%%%%%%%%
%%  Construcción del índice                                                 %%
%%%%%%%%%%%%%%%%%%%%%%%%%%%%%%%%%%%%%%%%%%%%%%%%%%%%%%%%%%%%%%%%%%%%%%%%%%%%%%%

\makeindex

%%%%%%%%%%%%%%%%%%%%%%%%%%%%%%%%%%%%%%%%%%%%%%%%%%%%%%%%%%%%%%%%%%%%%%%%%%%%%%
%%  Documento                                                              %%
%%%%%%%%%%%%%%%%%%%%%%%%%%%%%%%%%%%%%%%%%%%%%%%%%%%%%%%%%%%%%%%%%%%%%%%%%%%%%%

% \includeonly{Introduccion}

% \includexmp{licencia}

\begin{document}

\dominitoc

\begin{titlepage}
 \vspace*{2cm}
  \begin{center}
    {\huge \textbf{Álgebra constructiva en Haskell}}
  \end{center}
  \vspace{4cm}
  \begin{center}
    \leavevmode\includegraphics[totalheight=6cm]{sello.png}\\[3cm]
    {\normalsize Facultad de Matemáticas} \\
    {\normalsize Departamento de Ciencias de la Computación e Inteligencia Artificial}\\
    {\normalsize Trabajo Fin de Grado} \\
  \end{center}
  \begin{center}
    {\large \textbf{Ángela González Martín}}
  \end{center}
  \newpage
 
 \begin{flushright}
   \vspace*{5cm}
   \begin{minipage}{8.45cm}
      Agradecimientos
    \end{minipage}

      \vspace*{7.5mm}

  \end{flushright}
  \vspace*{\fill}

  \newpage


  
  % \begin{flushright}
  \begin{center}
   \vspace*{5cm}
    \begin{minipage}{14cm}
      El presente Trabajo Fin de Grado se ha realizado en el Departamento de
      Ciencias de la Computación e Inteligencia Artificial de la Universidad de
      Sevilla.

      \vspace*{7.5mm}

      Supervisado por
      % \vspace*{5mm}
    \end{minipage}\par
    María José Hidalgo Doblado
    % \end{flushright}
    \end{center}
  \vspace*{\fill}

  \newpage

  \vspace*{3cm}
  {\huge \textit{Abstract}}

  \vspace{2cm}
  Resumen en inglés
  
\end{titlepage}
\newpage

\noindent
Esta obra está bajo una licencia Reconocimiento--NoComercial--CompartirIgual
2.5 Spain de Creative Commons. 

\vspace*{4ex}

\begin{center}
\fbox{
\begin{minipage}{175mm}
\vspace*{1ex}
\noindent                     
{\bf Se permite:}
\begin{itemize}                         
\item copiar, distribuir y comunicar públicamente la obra
\item hacer obras derivadas
\end{itemize}
         
\noindent       
{\bf Bajo las condiciones siguientes:}

\begin{minipage}[t]{20mm}
  \vspace*{0mm}\hspace*{1em}
  \includegraphics[scale=0.35]{licencia/deed.png}
\end{minipage}%
\begin{minipage}[t]{15cm}
  \vspace*{0mm}
  {\bf Reconocimiento}. Debe reconocer los créditos de la  obra de la manera
  especificada por el autor.  
\end{minipage}
\\[1ex]

\begin{minipage}[t]{20mm}
  \vspace*{0mm}\hspace*{1em}
  \includegraphics[scale=0.35]{licencia/deed-eu.png}
\end{minipage}%
\begin{minipage}[t]{15cm}
  \vspace*{2mm}
  {\bf No comercial}. No puede utilizar esta obra para fines comerciales.
\end{minipage}
\\[1ex]

\begin{minipage}[t]{20mm}
  \vspace*{0mm}\hspace*{1em}
  \includegraphics[scale=0.35]{licencia/deed_002.png}
\end{minipage}%
\begin{minipage}[t]{15cm}
  \vspace*{-2mm}
  {\bf Compartir bajo la misma licencia}. Si altera o transforma esta obra, o
  genera una obra derivada, sólo puede distribuir la obra generada bajo una
  licencia idéntica a ésta. 
\end{minipage}

\begin{itemize}    
\item Al reutilizar o distribuir la obra, tiene que dejar bien claro los
  términos de la licencia de esta obra. 
\item Alguna de estas condiciones puede no aplicarse si se obtiene el permiso
 del titular de los derechos de autor.
\end{itemize}
\vspace*{1ex}
\end{minipage}}
\end{center}

\vspace*{4ex}

Esto es un resumen del texto legal (la licencia completa). Para ver una copia
de esta licencia, visite 
\href{http://creativecommons.org/licenses/by-nc-sa/2.5/es/}
     {\url{http://creativecommons.org/licenses/by-nc-sa/2.5/es/}}
o envie una carta a Creative Commons, 559 Nathan Abbott Way, Stanford,
California 94305, USA.

%%% Local Variables: 
%%% mode: latex
%%% TeX-master: t
%%% End: 

\newpage 

\tableofcontents
\newpage

\addcontentsline{toc}{chapter}{Introducción}\label{sec:introHas}
\chapter*{Introducción}

Este trabajo tiene como objetivo implementar conceptos de álgebra constructiva en Haskell, desde la noción de anillo hasta la de anillo coherente. El camino a seguir consiste en desarrollar los siguientes puntos:
\begin{itemize}
\item Familiarización con las funciones predefinidas en Haskell. Aquí aportaremos las funciones básicas de Haskell con el fin de facilitar la lectura del código, pues muchas de ellas son utilizadas a lo largo del trabajo. Algunas de ellas las podemos encontrar en la librería \href{http://hackage.haskell.org/package/base-4.11.1.0/docs/Data-List.html}{\texttt{Data.List}}  de Haskell.

\item Implementación de la teoría de anillos en Haskell. Introduciendo la noción de anillo, para posteriormente incluir los anillos conmutativos, dominio de integridad, cuerpos e ideales. Hemos definido los conceptos mediante programación funcional (\cite{Alonso-15b} y \cite{Hutton-16}) y teoría de tipos (\cite{Alonso-16a} y \cite{aprendehaskell}). De tal forma, que mediante las instancias podemos trabajar con cualquier anillo en concreto, como el anillo de los enteros, \mathbb{Z}. Para la elaboración de este capítulo se ha utilizado el material de la asignatura ``Estructuras Algebraicas'' del Departamento de Álgebra de la facultad de Matemáticas, ``Tema 3: Anillos'' (\cite{Algebra-15}). También se ha utilizado la tesis ``Constructive Algebra in Functional Programming and Type Theory'' de Anders Mörtberg (\cite{tesis}).

\item Librería de vectores y matrices sobre anillos en Haskell. Tomando como referencia la actual librería de matrices de la que Haskell dispone, \href{https://hackage.haskell.org/package/matrix-0.3.6.1/docs/Data-Matrix.html}{\texttt{Data.Matrix}}, junto con la tesis de Anders Mörtberg (\cite{tesis}). En este capítulo, comenzamos con las operaciones básicas entre matrices hasta poder aplicar el Método de Gauss, trabajando con vectores en forma de listas y matrices como listas de vectores sobre un anillo. Los conceptos sobre el método de Gauss-Jordan se encuentran en los apuntes de teoría ``Sistemas de ecuaciones lineales'' del Departamento de Álgebra (\cite{gauss}).

\item Finalmente, incorporamos la teoría sobre anillos coherentes en Haskell. El objetivo de este capítulo es poder resolver sistemas de ecuaciones de la forma $\,M\vec{X}=\vec{b}\,$ en anillos fuertemente discretos y coherentes. La primera sección del capítulo contiene la noción de anillo fuertemente discreto. En la siguiente sección, se define la noción de anillo coherente y se demuestran unas proposiciones que nos permiten implementar la resolución de sistemas de ecuaciones homogéneos. Al final del capítulo, trabajamos con anillos fuertemente discretos y coherentes, los cuáles nos permiten resolver sistemas de ecuaciones no homogéneos. Para el desarrollo de este capítulo, se ha hecho uso de las notas de teoría ``Coherent Ring'' de Thierry Coquand (\cite{coherent-10}) y para el código, se ha utilizado la tesis de Anders Mörtberg (\cite{tesis}).

\end{itemize}

Para poder llevar acabo este trabajo he necesitado utilizar varias herramientas informaticas, daré un breve resumen de cada una de ellas:\\

\begin{itemize}

\item \href{https://www.ubuntu.com/}{\textbf{Ubuntu}}\\
Es un sistema operativo de código abierto para ordenadores. Es una distribución de Linux basada en la arquitectura de Debian. Actualmente corre en ordenadores de escritorio y servidores, en arquitecturas Intel, AMD y ARM. Está orientado al usuario promedio, con un fuerte enfoque en la facilidad de uso y en mejorar la experiencia del usuario. Está compuesto de múltiple software normalmente distribuido bajo una licencia libre o de código abierto. Para desarrollar el trabajo hemos utilizado la versión Ubuntu 16.04 que se encuentra en la web oficial de Ubuntu, los pasos para su instalación están en el apéndice \ref{aped.A}. 

\item \href{https://www.latex-project.org/}{\textbf{LaTeX}}  \\
Es un sistema de composición de textos, orientado a la creación de documentos escritos que presenten una alta calidad tipográfica. Por sus características y posibilidades, se usa en la generación de artículos y libros científicos para incluir expresiones matemáticas, entre otros.  $\mathbf{L\!\!^{{}_{\scriptstyle A}} \!\!\!\!\!\;\; T\!_{\displaystyle E} \! X}$ está formado por un gran conjunto de macros de  $\mathbf{T\!_{\displaystyle E} \! X}$, con la intención de facilitar el uso del lenguaje de composición tipográfica ${\displaystyle \mathbf {T\!_{\displaystyle E}\!X} }$. $\mathbf{L\!\!^{{}_{\scriptstyle A}} \!\!\!\!\!\;\; T\!_{\displaystyle E} \! X}$ es software libre bajo licencia LPPL. 

\item \href{https://www.haskell.org/}{\textbf{Haskell}}\\
Es un lenguaje de programación puramente funcional con semánticas no estrictas y fuertemente tipado. Dentro de Haskell, encontramos \textbf{haskell literario}, con el cuál podemos escribir código en el lenguaje de programación de Haskell al mismo tiempo que podemos redactar texto con las utilidades que ofrece $\mathbf{L\!\!^{{}_{\scriptstyle A}} \!\!\!\!\!\;\; T\!_{\displaystyle E} \! X}$. Todo el trabajo está escrito en  \textbf{haskell literario}. Haskell 8.0.2 es la versión que se ha utilizado para realizar el trabajo.

\item \href{https://www.gnu.org/software/emacs/}{\textbf{Emacs}}\\
Es un editor de texto libre, extensible y personalizable. Emacs es el editor que he usado para programar en Haskell. Este dispone de muchas abreviaciones de teclado para facilitar su uso de una forma más rápida, las más utilizadas se pueden ver en el apéndice \ref{aped.B}.

\item \href{https://github.com/}{\textbf{GitHub}}\\
Es una plataforma de desarrollo colaborativo para alojar proyectos utilizando el sistema de control de versiones Git. Se utiliza principalmente para la creación de código fuente de programas. El código de los proyectos alojados en GitHub se almacena de forma pública, aunque utilizando una cuenta de pago, también permite tener repositorios privados. En mi \href{https://github.com/AngelaGM10}{perfil de GitHub} se encuentra el repositorio con el trabajo. Gracias a esta plataforma puedo compartir los avances con mi tutora. Esto es posible porque al crear un repositorio de GitHub, se pueden incluir colaboradores, que podrán añadir y modificar los elementos del repositorio. Para poder subir los archivos a GitHub podemos hacerlo desde la terminal de Ubuntu, los pasos a seguir están en el apéndice \ref{aped.C}.

\end{itemize}

\chapter{Programación funcional con Haskell}\label{sec:progfunHas}

En este capítulo se hace una breve introducción a la programación funcional en
Haskell suficiente para entender su aplicación en los siguientes
capítulos. Para una introducción más amplia se pueden consultar los apuntes de
la asignatura de Informática de 1º del Grado en Matemáticas
(\cite{Alonso-15b}). 

El contenido de este capítulo se encuentra en el módulo \texttt{PFH} 
\entrada{PFH}.

%%% Local Variables:
%%% mode: latex
%%% TeX-master: "TFG"
%%% End:


\chapter{Como empezar con Emacs en Ubuntu}

En este capítulo se hace una breve explicación de conceptos básicos para empezar a redactar un documento a LaTeX en Emacs y con Haskell a la vez, así como ir actualizando los archivos junto con la plataforma Github. Comenzaremos explicando como realizar la instalación de Ubuntu 16.04 en un PC con windows 10.\\

\section{Instalar Ubuntu 16.04 junto a windows 10}\label{aped.A}

Para realizar la instalación de Ubuntu junto a windows necesitaremos los siguientes programas:
\begin{itemize}
\item \href{https://rufus.akeo.ie/?locale=es_ES}{Rufus-2.17}

\item \href{https://www.ubuntu.com/download/desktop}{Ubuntu 16.04}

\end{itemize}
También necesitaremos un pen drive para usarlo de instalador.\\

\begin{itemize}
\item \textbf{Paso 1:}\\
Descargamos Ubuntu 16.04 y rufus-2.17 desde sus respectivas web (o enlaces dados anteriormente).\\
Necesitamos saber que tipo tiene nuestro disco duro, esto lo podemos ver haciendo click derecho sobre el icono de windows (abajo izquierda) y le damos a administrador de equipos -> administrador de discos, y nos aparecerá nuestro disco duro con todas sus subparticiones internas, en el general nos pondrá si es NTFS o MBR.\\

\item \textbf{Paso 2:}
\begin{itemize}
\item Conectamos el pen al PC y abrimos el programa rufus, el propio programa reconocerá el pen, sino en la pestaña de dispositivo marcamos el pen.

\item En Tipo de partición si nuestro disco es NTFS marcamos GPT para UEFI, en caso contrario uno de los otros dos MBR.

\item En la parte de opciones de formateo marcamos (aunque deben de venir marcadas): Formateo rápido

\item Crear disco de arranque con -> seleccionamos imagen ISO y hacemos click en el icono de la derecha para adjuntar la imagen ISO de Ubuntu que hemos descargado anteriormente.

\item Añadir etiquetas extendidas e iconos.

\item Seleccionamos empezar.\\
\end{itemize}

\item \textbf{Paso 3:}\\
Dejamos el pen conectado al PC y reiniciamos el ordenador, al reiniciar justo antes de que cargue pulsamos F2 (o F1 según el PC) para acceder a la bios del PC y aqui nos vamos a la zona de arranque de cada sistema (esto cada bios es diferente). Tenemos que colocar el pen en la primera posición, que en esta debe estar windows, de esta forma iniciamos con el pen y comenzamos a instalar Ubuntu. Seguimos los pasos que el programa de instalación nos indica, solo tenemos que marcar España cuando nos lo pida y dar el espacio que queramos a Ubuntu (con unos 40 GB es suficiente), el propio programa se encarga de hacer la partición del disco duro.\\

\item \textbf{Paso 4:}\\
Una vez instalado Ubuntu, nos vamos al icono naranjita que se llama software de Ubuntu y actualizamos.\\

Tras realizar todos estos pasos, cuando iniciemos el PC nos debe dar a elegir entre iniciar con Ubuntu o con Windows 10. Recomiendo buscar en youtube un video tutorial de como instalar Ubuntu junto a windows 10.\\
\end{itemize}

\section{Iniciar un Capítulo}

El editor de Emacs podemos descargarlo de su página web, la instalación es muy sencilla y guiada. Del mismo modo, instalamos Haskell 8.0.2. Una vez tenemos ambos programas instalados, procedemos a iniciar un capítulo en Haskell literario.
\begin{itemize}

\item \textbf{Paso 1:}\\
Abrimos el directorio desde Emacs con \fbox{Ctrl+x+d} y accedemos a la carpeta de texto para crear el archivo nuevo .tex sin espacios. Aquí escribiremos la introducción a nuestro capítulo y secciones. Desde este fichero .tex se importarán los módulos.

\item \textbf{Paso 2:}\\
Hacemos lo mismo pero en la carpeta código y guardamos el archivo con la abreviatura que hemos usado en el .tex, el archivo lo guardamos como .lhs para tener ahí el código necesario de Haskell. Estos serán los módulos donde desarrollaremos las secciones con Haskell literario, es decir, escribiremos a Latex y a su vez tendremos el código de Haskell.\\

\item \textbf{Paso 3:}\\
Al acabar el capitulo hay que añadir el capítulo al resto del trabajo. Para ello, nos vamos al archivo que contiene todo el trabajo que en nuestro caso se llama 'TFG.tex' , importante coger el de la extensión .tex, nos vamos a la zona donde incluimos los capitulos y usamos el comando de LaTeX con el nombre que le dimos en la carpeta de texto:\\

\fbox{\ include\{'nombre sin el .tex'\}}\\

\end{itemize}

\section{Abreviaciones de Emacs}\label{aped.B}\\

La tecla ctrl se denominara C y la tecla alt M, son las teclas mas utilizadas, pues bien ahora explicamos los atajos más importantes y seguiremos la misma nomenclatura de la guía para las teclas:\\

ctrl es llamada C y alt M.\\

Para abrir o crear un archivo:\\
\fbox{C + x + C + f}\\

Para guardar un archivo:\\
\fbox{C + x + C + s}\\

Para guardar un archivo (guardar como):\\
\fbox{C + x + C + w}\\

Si abriste mas de un archivo puedes recorrerlos diferentes buffers con\\
\fbox{C + x + ← o →}\\

Emacs se divide y maneja en buffers y puedes ver varios buffers a la vez (los buffers son como una especie de ventanas).\\

Para tener 2 buffers horizontales:\\
\fbox{C + x + 2}\\

Para tener 2 buffers verticales:\\
\fbox{C + x + 3}\\

Para cambiar el puntero a otro buffer:\\
\fbox{C + x + o}\\

Para tener un solo buffer:\\
\fbox{C + x + 1}\\

Para cerrar un buffer:\\
\fbox{C + x + k}\\

Si por ejemplo nos equivocamos en un atajo podemos cancelarlo con:\\
\fbox{C + g}\\

Para cerrar emacs basta con:\\
\fbox{C + x + C + C}\\

Para suspenderlo:\\
\fbox{C + z}\\

\section{Push and Pull de Github con Emacs}\label{aped.C}

Vamos a mostrar como subir y actualizar los archivos en la web de Github desde la Consola (o Terminal). Lo primero que debemos hacer es abrir la Consola:\\

\fbox{Ctrl+Alt+T}\\

Escribimos los siguientes comandos en orden para subir los archivos:\\

\fbox{cd 'directorio de la carpeta en la que se encuentran las subcarpetas de código y texto'} ejemplo: cd Escritorio/AlgebraConEnHaskell/ \\

\fbox{git add .} (de esta forma seleccionamos todo)\\

\fbox{git commit -m 'nombre del cambio que hemos hecho'}\\

\fbox{git push origin master}\\

A continuación, nos pedirá el nombre de usuario de GitHub, lo ponemos y seguidamente nos pedirá la contraseña, al ponerla se subirán los archivos a la plataforma de GitHub.


Para descargar los archivos hacemos lo mismo cambiando el último paso por:\\

\fbox{git pull origin master}\\


%%% Local Variables:
%%% mode: latex
%%% TeX-master: "TFG"
%%% End:

\chapter{Teoría de anillos en Haskell}\label{sec:anillosHas}

En este capítulo se muestra cómo definir la teoría de anillos en Haskell. Los anillos se pueden definir de forma compacta en grupos y monoides, pero daremos unas series de definiciones que los definen de forma más rigurosa.\\

Comenzamos dando las primeras definiciones y propiedades básicas que tiene un anillo en el módulo \texttt{TAH} 
\entrada{TAH}

Para describir los anillos commutativos necesitamos un nuevo módulo \texttt{TAHCommutative}\entrada{TAHCommutative}

Dentro de este último módulo podemos definir el dominio de integridad, pues se necesita que el anillo sea conmutativo, usaremos el módulo \texttt{TAHIntegralDomain}\entrada{TAHIntegralDomain}

Utilizaremos un nuevo módulo para dar la definición de cuerpo que se encuentra en \texttt{TAHCuerpo}\entrada{TAHCuerpo}

\section{Ideales}
El concepto de ideales es muy importante en el álgebra conmutativa. Son
generalizaciones de muchos conceptos de los enteros. Dado que solo consideramos anillos conmutativos, la propiedad multiplicativa de izquierda y derecha son la misma. Veamos su implementación en 

\section{Anillos fuertemente discretos}
Esta sección considerará algunos anillos que son especialmente relevantes para la construcción matemática. Empezamos con el módulo \texttt{TAHStronlyDiscrete}\entrada{TAHStronglyDiscrete}

%%% Local Variables:
%%% mode: latex
%%% TeX-master: "TFG"
%%% End:


\chapter{Vectores y Matrices}\label{sec:matrixHas}
En este capítulo implementaremos los vectores y matrices como listas para poder usarlas en los módulos siguientes. En Haskell existen ya estas librerías pero solo podemos usarlas si nos restringimos al tipo numérico. Por ello, crearemos este módulo con el fin de generalizar los conceptos para anillos sobre cualquier tipo $a$.

\section{Vectores}
Antes de introducir las matrices tenemos que especificar ciertas operaciones y funciones sobre los vectores, pues las filas de las matrices serán vectores, así tendremos la matriz como lista de listas. Damos los conceptos necesarios en el módulo \texttt{TAHVector}\entrada{TAHVector} 

\section{Matrices}
Una vez dadas las funciones y operaciones de los vectores podemos comenzar a implementar las matrices, notesé que cada fila o columna de una matriz puede verse como un vector. Nuestra idea es dar las matrices como una lista de vectores de forma que cada vector será una fila de la matriz. El objetivo de esta sección es implementar el método de Gauss-Jordan con el fin de poder resolver sistemas de la forma $Ax=b$. Lo vemos en el módulo \texttt{TAHMatrix}\entrada{TAHMatrix} 

\chapter{Anillos Coherentes}\label{sec:coherentHas}
Todos los anillos en este capítulo son dominios de integridad. Uno de los principales objetivos de las siguientes secciones será ver cuando un anillo es coherente. Es decir, anillos en los que es posible resolver sistemas de ecuaciones. Antes de comenzar, introduciremos la noción de anillo fuertemente discreto.
\section{Anillos Fuertemente Discretos}
En esta breve sección mostraremos la noción de anillo discreto y fuertemente discreto, lo vemos en el módulo \texttt{TAHStronglyDiscrete}\entrada{TAHStronglyDiscrete} 

\section{Anillos Coherentes}
En esta sección, nos ayudaremos del módulo de vectores y matrices creado en el capítulo anterior para construir, en Haskell, la noción de anillo coherente con el objetivo de poder resolver sistemas de ecuaciones. Lo vemos en el módulo \texttt{TAHCoherent}\entrada{TAHCoherent} 
%%% Local Variables:
%%% mode: latex
%%% TeX-master: "TFG"
%%% End:

%%%%%%%%%%%%%%%%%%%%%%%%%%%%%%%%%%%%%%%%%%%%%%%%%%%%%%%%%%%%%%%%%%%%%%%%%%%%%%% 
%%  Apéndice                                                         %%
%%%%%%%%%%%%%%%%%%%%%%%%%%%%%%%%%%%%%%%%%%%%%%%%%%%%%%%%%%%%%%%%%%%%%%%%%%%%%%%

\appendix
\chapter{Más cosas}\label{aped.A}
Aún faltan cosas por decir.

\chapter{Y más cosas aún}\label{aped.B}
Y más cosas aún.

%%%%%%%%%%%%%%%%%%%%%%%%%%%%%%%%%%%%%%%%%%%%%%%%%%%%%%%%%%%%%%%%%%%%%%%%%%%%%%% 
%%  Bibliografía                                                            %%
%%%%%%%%%%%%%%%%%%%%%%%%%%%%%%%%%%%%%%%%%%%%%%%%%%%%%%%%%%%%%%%%%%%%%%%%%%%%%%%

\addcontentsline{toc}{chapter}{Bibliografía}
\bibliographystyle{abbrv}
\bibliography{TFG}

%%%%%%%%%%%%%%%%%%%%%%%%%%%%%%%%%%%%%%%%%%%%%%%%%%%%%%%%%%%%%%%%%%%%%%%%%%%%%%%
%%  Índice                                                                  %%
%%%%%%%%%%%%%%%%%%%%%%%%%%%%%%%%%%%%%%%%%%%%%%%%%%%%%%%%%%%%%%%%%%%%%%%%%%%%%%%

\addcontentsline{toc}{chapter}{Indice de definiciones}

\printindex

%%%%%%%%%%%%%%%%%%%%%%%%%%%%%%%%%%%%%%%%%%%%%%%%%%%%%%%%%%%%%%%%%%%%%%%%%%%%%%%
%% § Pendientes                                                              %%
%%%%%%%%%%%%%%%%%%%%%%%%%%%%%%%%%%%%%%%%%%%%%%%%%%%%%%%%%%%%%%%%%%%%%%%%%%%%%%%

%\todototoc
%\listoftodos

\end{document}



%%% Local Variables:
%%% mode: latex
%%% TeX-master: t
%%% End:
